\subsection{\texorpdfstring{Integrability criteria in \(\mathbf{R}^{n}\) and \(\mathbf{Z}^{n}\)}{Integrability criteria in \textbackslash mathbf\{R\}\^{}\{n\} and \textbackslash mathbf\{Z\}\^{}\{n\}}}

PROPOSITION 3. --- Let N be a norm on \(\mathbf{R}^{n}\) , \(r\) a real number and \(p \in [1, +\infty[\) .

\begin{enumerate}
\def\labelenumi{\alph{enumi})}
\tightlist
\item
  The function \((1 + \mathrm{N})^{r}\) belongs to \(\mathcal{L}^{p}(\mathbf{R}^{n}, \mu)\) if and only if \(rp < -n\);
\end{enumerate}

a') The restriction to \(\mathbf{Z}^{n}\) of the function \((1+N)^{r}\) belongs to \(\ell^{p}(\mathbf{Z}^{n})\) if and only if \(rp < -n\);

\begin{enumerate}
\def\labelenumi{\alph{enumi})}
\setcounter{enumi}{1}
\tightlist
\item
  Let V be a bounded measurable neighborhood of 0 in \(\mathbf{R}^{n}\) . The function \(\mathbf{N}^{r}\) belongs to \(\mathcal{L}^{p}(\mathbf{R}^{n}-\mathbf{V},\mu)\) if and only if \(rp<-n\) .
\end{enumerate}

By virtue of the equivalence of norms on \(\mathbf{R}^{n}\) (EVT I, p. 14, th. 2 and TG, IX, p. 32, prop. 8), one may assume that the norm N is given by \(\mathrm{N}(x)=\sup|x_{i}|\) for \(x=(x_{i})\in\mathbf{R}^{n}\) . Let B denote the unit ball of \(\mathbf{R}^{n}\) for this norm.

Let V be a bounded measurable neighborhood of 0 in \(\mathbf{R}^n\) . The function \(N^r\) is continuous and bounded on \((\mathrm{B - V})\cup (\mathrm{V - B})\) , which shows that the assertion \(b)\) is valid if and only if it is when \(V = B\) , which we will assume from now on. The function \((1 + N)^r\) being continuous and bounded on B, and satisfying the inequality \(N^r\leqslant (1 + N)^r\leqslant 2^r N^r\) on \(\mathbf{R}^n -\mathbf{B}\) one observes that the assertions \(a)\) and \(b)\) are equivalent.

Let us prove b) when V = B. We may assume that p = 1. The case where \(r > 0\) being elementary, suppose that \(r \leqslant 0\) . For every integer \(j \geqslant 1\) the integrable set

\[
\mathrm{C} _ {j} = \{x \in \mathbf {R} ^ {n} \mid 2 ^ {j - 1} \leqslant \mathrm{N} (x) <   2 ^ {j} \}
\]

has measure \(2^{nj}(2^n - 1)\) . The sets \((\mathrm{C}_j)_{j \geqslant 1}\) form a partition of \(\mathbf{R}^n - \mathbf{B}\) . By INT, V, p. 4, § 1, no. 1, cor., we therefore have

\[
\begin{array}{c} \int_ {\mathbf {R} ^ {n} - \mathrm{B}} ^ {*} \mathrm{N} ^ {r} d \mu = \sum_ {j \geqslant 1} \int_ {\mathrm{C} _ {j}} ^ {*} \mathrm{N} ^ {r} d \mu \\ \leqslant \sum_ {j \geqslant 1} 2 ^ {n + n j} 2 ^ {r (j - 1)} = 2 ^ {n - r} \sum_ {j \geqslant 1} 2 ^ {(n + r) j} \end{array}
\]

which is finite if \(r < -n\) . On the other hand, we have (loc. cit.)

\[
\int_ {\mathbf {R} ^ {n} - \mathrm{B}} ^ {*} \mathrm{N} ^ {r} d \mu \geqslant \sum_ {j \geqslant 1} 2 ^ {r j} 2 ^ {n j} (2 ^ {n} - 1) = (2 ^ {n} - 1) \sum_ {j \geqslant 1} 2 ^ {(r + n) j}
\]

which is infinite if \(r \geqslant -n\) .

One proves \(a'\) ) in a similar manner by considering the sets \(\mathrm{C}_{j} \cap \mathbf{Z}^{n}\) which cover \(\mathbf{Z}^{n} - \{0\}\) and satisfy

\[
\operatorname{Card} \left(\mathrm{C} _ {j} \cap \mathbf {Z} ^ {n}\right) = (2 ^ {j + 1} - 1) ^ {n} - (2 ^ {j} - 1) ^ {n},
\]

which belongs to the interval \([(1-2^{-n})(2^{j+1}-1)^{n},2^{n(j+1)}]\) .

